\documentclass[11pt]{article}
\usepackage[margin=1in]{geometry}
\usepackage{array}
\usepackage{enumitem}
\usepackage{xcolor}
\usepackage{hyperref}
\usepackage{listings}

\hypersetup{colorlinks=true, linkcolor=blue, urlcolor=blue}
\setlength{\parindent}{0pt}
\setlength{\parskip}{6pt}
\setlist[itemize]{topsep=3pt,itemsep=2pt,leftmargin=*}
\setlist[enumerate]{topsep=3pt,itemsep=4pt,leftmargin=*}

\lstset{
  basicstyle=\ttfamily\small,
  breaklines=true,
  frame=single,
  columns=fullflexible,
  keepspaces=true,
  showstringspaces=false,
  backgroundcolor=\color{gray!6}
}

\definecolor{predictcolor}{RGB}{255,242,200}
\definecolor{discusscolor}{RGB}{214,234,255}
\definecolor{funfactcolor}{RGB}{222,247,222}
\definecolor{debatecolor}{RGB}{255,221,221}

\newcommand{\calloutbox}[3]{%
  \vspace{4pt}\noindent
  \colorbox{#1}{\parbox{0.96\linewidth}{\textbf{#2}}}\\[2pt]
  \noindent\fbox{\begin{minipage}{0.94\linewidth}\vspace{2pt}#3\vspace{2pt}\end{minipage}}
  \vspace{6pt}
}
\newcommand{\predictbox}[1]{\calloutbox{predictcolor}{PREDICTION TIME}{#1}}
\newcommand{\discussbox}[1]{\calloutbox{discusscolor}{HUDDLE UP}{#1}}
\newcommand{\funfact}[1]{\calloutbox{funfactcolor}{FUN FACT}{#1}}
\newcommand{\debatebox}[1]{\calloutbox{debatecolor}{DEBATE CORNER}{#1}}
\newcommand{\claimbox}[1]{\calloutbox{funfactcolor}{CLAIM, EVIDENCE, CONTEXT}{#1}}
\newcommand{\badge}[1]{$[\ ]$ \textbf{#1}}

\newcommand{\writebox}[1]{
  \vspace{3pt}
  \noindent\fbox{
    \begin{minipage}{0.96\linewidth}
      \textbf{#1}\\[12pt]
      \rule{\linewidth}{0.4pt}\\[14pt]
      \rule{\linewidth}{0.4pt}\\[14pt]
      \rule{\linewidth}{0.4pt}\\[14pt]
    \end{minipage}
  }
  \vspace{6pt}
}

\newcommand{\shortwritebox}[1]{
  \vspace{3pt}
  \noindent\fbox{
    \begin{minipage}{0.96\linewidth}
      \textbf{#1}\\[12pt]
      \rule{\linewidth}{0.4pt}\\[14pt]
    \end{minipage}
  }
  \vspace{6pt}
}

\begin{document}

\begin{center}
  {\LARGE MA 213 Week 3: Lab 3}\\[3pt]
  {\large From Data to Visual Evidence}
\end{center}

\begin{enumerate}
  \item \textbf{Your name:} \underline{\hspace{0.3\textwidth}} \hfill \textbf{BU ID:} \underline{\hspace{0.25\textwidth}}
  \item \textbf{Partner's name:} \underline{\hspace{0.3\textwidth}} \hfill \textbf{BU ID:} \underline{\hspace{0.25\textwidth}}
\end{enumerate}

\textbf{Big idea:} A good plot is not just a picture. It is evidence. In this lab, you will move from raw data to clear summaries, plots, and interpretations about air quality and Titanic passenger survival.
\textbf{Do not use GenAI tools for this lab.} The point is to practice your own analysis work, coding skills, and statistical reasoning.

\textbf{Files used:} \texttt{airquality.csv} and \texttt{Titanic.csv}

\textbf{Skills practiced:} \texttt{read.csv()}, \texttt{filter()}, \texttt{mutate()}, \texttt{group\_by()}, \texttt{summarize()}, \texttt{count()}, and \texttt{ggplot()}.

\textbf{Your mission:} Work with your partner to decide which variables are categorical or numerical, choose appropriate visualizations, and write claims that are supported by plots or tables.

\section*{Icebreaker}

Before opening R: introduce yourself to your partner. Then answer: What is one environmental condition that affects how you feel during a day, such as heat, wind, rain, or air quality?

\shortwritebox{My answer and my partner's answer:}

\section*{Getting Started: Open the Evidence Files}

Load the packages and import both data sets. The first column in each CSV is a row index, so we use \texttt{row.names = 1}.

\begin{lstlisting}[language=R]
library(dplyr)
library(ggplot2)
library(tidyr)

air <- read.csv("airquality.csv", row.names = 1)
titanic <- read.csv("Titanic.csv", row.names = 1)

head(air)
head(titanic)

str(air)
str(titanic)
\end{lstlisting}

\section*{Question 1. What kind of data are we looking at?}

\discussbox{Before calculating anything, discuss with your partner: What does one row represent in each data set? Which variables are numerical? Which variables are categorical?}

\discussbox{What patterns or relationships might we expect to find in each data set?}

\writebox{Write one possible research question for each data set.}

\funfact{The \texttt{airquality} data contains daily air quality measurements from New York, May through September 1973. The Titanic data records passenger characteristics and whether passengers survived.}

\section*{Question 2. Does solar radiation vary by month?}

Use the air quality data to compare average solar radiation across months.

\predictbox{Before running code, which month do you think will have the highest average solar radiation? Why?}

\begin{lstlisting}[language=R]
solar_by_month <- air %>%
  group_by(________) %>%
  summarize(solar_r_avg = mean(________, na.rm = TRUE))

solar_by_month

ggplot(solar_by_month, aes(x = factor(________), y = solar_r_avg)) +
  geom_col() +
  labs(
    x = "Month",
    y = "Average solar radiation",
    title = "Average solar radiation by month"
  )
\end{lstlisting}

\writebox{Which month has the highest average solar radiation? Was your prediction correct?}

\section*{Question 3. How are temperature and solar radiation related?}

A scatter plot is useful when both variables are numerical. Create a plot comparing \texttt{Temp} and \texttt{Solar.R}.

\begin{lstlisting}[language=R]
air_clean <- air %>%
  filter(!is.na(Solar.R), !is.na(Temp))

ggplot(air_clean, aes(x = Temp, y = Solar.R)) +
  geom_point() +
  labs(
    x = "Temperature (F)",
    y = "Solar radiation",
    title = "Solar radiation and temperature"
  )
\end{lstlisting}

\discussbox{Does the scatter plot suggest a positive association, negative association, or no clear association? Are there any points that seem unusual?}

\writebox{Describe the relationship between temperature and solar radiation in context.}

\section*{Question 4. Can we categorize air quality conditions?}

Look at how the code below uses \texttt{mutate()} to create simplified Beaufort-style categories (recall the Beaufort wind scale from Tutorial 2). What is each part of the code doing? How are the new variables \texttt{ozone\_level} and \texttt{wind\_level} created?

\begin{lstlisting}[language=R]
air_categories <- air %>%
  filter(!is.na(Ozone), !is.na(Wind)) %>%
  mutate(
    ozone_level = ifelse(Ozone > median(Ozone), "High", "Low"),
    wind_level = case_when(
      Wind < 1 ~ "Calm",
      Wind < 4 ~ "Light air",
      Wind < 7 ~ "Light breeze",
      Wind < 12 ~ "Gentle breeze",
      Wind < 18 ~ "Moderate breeze",
      Wind < 24 ~ "Fresh breeze",
      TRUE ~ "Strong breeze"
    )
  )

ozone_wind_table <- air_categories %>%
  count(wind_level, ozone_level)

ozone_wind_table
\end{lstlisting}

\writebox{What does the \texttt{filter()} line do here?}

\writebox{The \texttt{ifelse()} function takes three arguments: a condition, a value if TRUE, and a value if FALSE. In words, what rule does \texttt{ifelse()} use to assign \texttt{ozone\_level}?}

\writebox{The \texttt{case\_when()} function checks a list of conditions in order and uses the value tied to the first one that is TRUE. Why does \texttt{case\_when()} make more sense than \texttt{ifelse()} for this variable?}

\debatebox{The original \texttt{Wind} variable is numerical. After you create \texttt{wind\_level}, should it be treated as numerical or categorical? Make your case.}

\section*{Question 5. Did Titanic survival differ by class?}

Now investigate two categorical variables: passenger class and survival.

\predictbox{Before making the table, which passenger class do you think had the highest survival count? Which class do you think had the highest survival rate? These may not be the same question.}

\begin{lstlisting}[language=R]
class_survival_table <- titanic %>%
  count(________, ________)

class_survival_table

ggplot(class_survival_table, aes(x = Class, y = n, fill = Survived)) +
  geom_col(position = "dodge") +
  labs(
    x = "Passenger class",
    y = "Number of passengers",
    title = "Titanic survival counts by class"
  )
\end{lstlisting}

\writebox{Which class had the highest survival count? What general pattern do you see between passenger class and survival?}

\section*{Question 6. Counts or proportions?}

Counts can be helpful, but they can also hide important patterns when groups have different sizes. Use a filled bar chart to compare survival proportions by class.

\begin{lstlisting}[language=R]
ggplot(class_survival_table, aes(x = Class, y = n, fill = Survived)) +
  geom_col(position = "fill") +
  labs(
    x = "Passenger class",
    y = "Proportion within class",
    title = "Titanic survival proportions by class"
  )
\end{lstlisting}

\claimbox{In the Titanic data, passenger class and survival \underline{\hspace{5cm}} (do/do not) appear to be related. My evidence is \underline{\hspace{5cm}}.}

\section*{Optional Challenge. Did survival differ by sex or age?}

Using the same approach as above, investigate whether survival differed by \texttt{Sex} and by \texttt{Age}. 

\writebox{Which sex had the highest survival count? Which had the highest survival rate?}

\writebox{Which age group had the highest survival rate? Was the difference between children and adults large or small?}

\discussbox{Does the pattern you see for \texttt{Sex} resemble the pattern you saw for \texttt{Class}? What about the pattern for \texttt{Age}?}

\claimbox{In the Titanic data, \underline{\hspace{3cm}} appears to be more strongly related to survival than \underline{\hspace{3cm}}. My evidence is \underline{\hspace{5cm}}.}

\section*{Extra Challenge. Is there an interaction between sex and class?}

So far you have looked at survival by class alone and by sex alone. But the effect of class on survival might not be the same for both sexes -- this is called an \textbf{interaction}. A mosaic plot can show three categorical variables at once: the width of each column reflects the number of passengers in that group, and the height of the colored blocks within a column reflects the proportion who survived.

\predictbox{Do you think the ``women and children first'' policy applied equally across all three passenger classes? Why or why not?}

\begin{lstlisting}[language=R]
sex_class_survival_table <- table(titanic$Sex, titanic$Class, titanic$Survived)
mosaicplot(sex_class_survival_table, color=TRUE)
\end{lstlisting}

\discussbox{Look at the survival split within each Class-by-Sex combination. In which combination is the survival rate highest? In which is it lowest? Does the class pattern look the same for men as it does for women?}

\claimbox{The relationship between class and survival \underline{\hspace{3cm}} (does / does not) appear to depend on sex. I believe this because \underline{\hspace{5cm}}.}


\end{document}
